\documentclass{article}

\usepackage[a4paper, margin=0.5in]{geometry}
\usepackage{titling}
\usepackage{hyperref}

\setlength{\droptitle}{-5em}

\input{../macro/pluverse-preamble.tex}

\begin{document}

\title{Anonymized IDs}
\author{Hongxu Xu}

\maketitle

To anonymize IDs, you should add this line before using \texttt{\textbackslash anonymizedId}:
\begin{verbatim}
    \def\useAnonymizedId{}
    We have ids: \anonymizedId{1234}, \anonymizedId{5678}.

    We can reference them again:
        \anonymizedId{1234}, \anonymizedId{5678}.

    The anonymized ids should be consistent.
\end{verbatim}

\begin{tcolorbox}
\def\useAnonymizedId{}
We have ids: \anonymizedId{1234}, \anonymizedId{5678}.

We can reference them again:
    \anonymizedId{1234}, \anonymizedId{5678}.

The anonymized ids should be consistent.
\end{tcolorbox}

You can also set it according to the anonymous defined in acmart class:
\begin{verbatim}
\makeatletter
\if@ACM@anonymous
    \def\useAnonymizedId{}
\fi
\makeatother
\end{verbatim}

You can turn off the anonymization at any time by undefining:
\begin{verbatim}
    \undef\useAnonymizedId
    We have ids: \anonymizedId{1234}, \anonymizedId{5678}.
\end{verbatim}

\begin{tcolorbox}
\undef\useAnonymizedId
We have ids: \anonymizedId{1234}, \anonymizedId{5678}.
\end{tcolorbox}

Of course, you can wrap the number with hyperlinks or other commands as needed.
\begin{verbatim}
    \newcommand{\issueId}[1]{%
    \href{https://github.com/llvm/llvm-project/issues/#1}{#1}%
    }

    \newcommand{\anonymizedIssueId}[1]{%

    \ifx\useAnonymizedId\undefined%
        \issueId{#1}%
    \else%
        \myGetOrAssignID{#1}%
    \fi}

    \def\useAnonymizedId{}
    We have ids:
        \anonymizedIssueId{1234}, \anonymizedIssueId{5678}.
    Ref again:
        \anonymizedIssueId{1234}, \anonymizedIssueId{5678}.

    \undef\useAnonymizedId
    We have ids:
        \anonymizedIssueId{1234}, \anonymizedIssueId{5678}.
    Ref again:
        \anonymizedIssueId{1234}, \anonymizedIssueId{5678}.
\end{verbatim}

\newcommand{\issueId}[1]{%
    \href{https://github.com/llvm/llvm-project/issues/#1}{#1}%
}
\newcommand{\anonymizedIssueId}[1]{\ifx\useAnonymizedId\undefined%
  \issueId{#1}%
\else%
  \myGetOrAssignID{#1}%
\fi}

\begin{tcolorbox}
\def\useAnonymizedId{}
We have ids: \anonymizedIssueId{1234}, \anonymizedIssueId{5678}.
Ref again: \anonymizedIssueId{1234}, \anonymizedIssueId{5678}.

\undef\useAnonymizedId
We have ids: \anonymizedIssueId{1234}, \anonymizedIssueId{5678}.
Ref again: \anonymizedIssueId{1234}, \anonymizedIssueId{5678}.
\end{tcolorbox}

\end{document}
